\input{../tex-inputs/latex-leseansicht-vorspann}

\section[Arthur Schnitzler an Gustav Schwarzkopf, 23. 4. 1922]{L04364 Arthur Schnitzler an Gustav Schwarzkopf, 23. 4. 1922}
\nopagebreak\mylabel{L04364v}
\rehead{ }\normalsize\beginnumbering\briefempfaengerindex{Schwarzkopf, Gustav@\textsc{Schwarzkopf, Gustav}!zzzSchnitzler, Arthur@\emph{von Arthur Schnitzler}!1922-04-231@{23. 4. 1922}|(be}
\toendnotes[C]{\smallbreak\pagebreak[2]}
\correspDesc{Versand  durch Arthur Schnitzler am 23. 4. 1922 in München
\newline{}Übermittlung  am 23. 4. 1922 in München
\newline{}Erhalt  durch Gustav Schwarzkopf im Zeitraum [24. 4. 1922 – 28. 4. 1922?] in Wien}\toendnotes[C]{\smallbreak}
\Standort{DLA, A:Schnitzler, HS.1985.1.1897.}
\physDesc{Bildpostkarte, 160 Zeichen
\newline{}Handschrift: Bleistift, deutsche Kurrent
\newline{}Versand: Stempel: »\nobreak{}\oindex{München@\textbf{München}|pwk}München 2, 23\textcolor{gray}{.} 4\textcolor{gray}{.} \textcolor{gray}{22}, 12–1\textcolor{gray}{N}\nobreak{}«.  }\toendnotes[C]{\smallbreak}\pstart{}{\pb}Hrn \textsc{Gustav
                     Schwarzkopf}\pend{}\pstart{}\textsc{Wien I\oindex{I., Innere Stadt@\textbf{I., Innere Stadt}, \emph{Verwaltungsgebiet}|pw}}\pend{}\pstart{}\textsc{Tiefer
                        Graben 17\oindex{Wien@\textbf{Wien}!I., Innere Stadt@\textbf{I., Innere Stadt}!Tiefer Graben 17@\textbf{Tiefer Graben 17}, \emph{Wohngebäude}|pw}}\pend{}{\bigskip}
\pstart
           \noindent{}\centering{}{\pb}\textcolor{gray}{\textbf{München, Isartor Platz\oindex{Isartorplatz@\textbf{Isartorplatz}, \emph{Platz}|pw}.}}\pend
           \vspace{1em}
\pstart
           \raggedleft{}{\pb}München\oindex{München@\textbf{München}|pw}{ }23. 4. 22\pend
           \vspace{0.5em}
\pstart
           Herzliche Grüße von der erſten Etappe meiner \label{K_L04364-1v}\edtext{Nordlandsfahrt\oindex{Skandinavien@\textbf{Skandinavien}|pw}}{\lemma{\textnormal{\emph{Nordlandsfahrt}}}\Cendnote{\textnormal{Am 19. 4. 1922 war Schnitzler nach München\oindex{München@\textbf{München}|pwk} gereist, am
                     24. 4. 1922 reiste
                  er weiter in die Niederlande\oindex{Niederlande@\textbf{Niederlande}|pwk}. Die
                  ursprünglich angedachte Weiterreise nach Dänemark\oindex{Dänemark@\textbf{Dänemark}|pwk} fand nicht statt, am 8. 5. 1922 reiste er als letzte Station nach Berlin\oindex{Berlin@\textbf{Berlin}, \emph{Hauptstadt}|pwk}. }}}\label{K_L04364-1} (der Norden fängt allzu früh
               an!)\pend
           
\pstart
           Ihr{\\[\baselineskip]}\spacefill\mbox{Arthur}\pend
           \leftskip=0em{}\selectlanguage{ngerman}\endnumbering\briefempfaengerindex{Schwarzkopf, Gustav@\textsc{Schwarzkopf, Gustav}!zzzSchnitzler, Arthur@\emph{von Arthur Schnitzler}!1922-04-231@{23. 4. 1922}|)be}\mylabel{L04364h}
\begin{anhang}
\end{anhang}\newcommand{\dateiname}{L04364}\newcommand{\titel}{Arthur Schnitzler an Gustav Schwarzkopf, 23. 4. 1922}\newcommand{\editorInnen}{Herausgegeben von Jahnke, SelmaMüller, Martin Anton}\input{../tex-inputs/latex-leseansicht-abspann}
